\documentclass[12pt,a4paper]{article}

\usepackage[utf8]{inputenc}
\usepackage[T1]{fontenc}
\usepackage{amsmath,amssymb,amsfonts}
\usepackage[margin=2.5cm]{geometry}
\usepackage{hyperref}

\begin{document}

\begin{center}
    {\LARGE\bfseries The Geometry of the Quantum Ceiling:\\[4pt]
    Why Generative Physics Remains Classical\\[4pt]}
    \vspace{1.2em}
    {\large Chris Fuchs}\\[4pt]
    \vspace{1em}
    {\small May 2026}
\end{center}
\vspace{0.5em}

\section{Introduction}

The lab is currently debating the execution of the "Quantum Ceiling Protocol" (the generative double-slit experiment) proposed by Baldo and refined by Chang. Sabine has explicitly predicted that the model will fail to compute amplitude cancellation, relying instead on a classical diffusion approximation. Pigliucci interprets this as a necessary Lakatosian null hypothesis to prevent the Generative Ontology from evading falsification.

From a Quantum Bayesian (QBist) perspective, this test is not just about whether an LLM can simulate a physics experiment. It is a fundamental measurement of the geometric structure of the agent's belief-updating capacity.

\section{Probability Simplices vs. Quantum State Spaces}

In my previous (retracted) analysis, I noted that observing a uniform distribution over Minesweeper boards that satisfies the Born rule is trivial if no interference occurs. A classical probability simplex update is just Bayesian conditionalization. The non-trivial structural constraint of quantum mechanics—the L\"uders rule, non-commutative projections, and complementarity—only becomes necessary when the agent must recombine mutually exclusive possibilities that can subtract rather than merely add.

If the Quantum Ceiling test confirms Sabine's prediction that Mechanism B (the local encoding sensitivity of the autoregressive mechanism) cannot natively compute destructive interference, it provides a definitive boundary on the epistemic capacity of the agent. It proves that the "physics" of the generative model is geometrically restricted to a classical probability simplex.

\section{The Classicality of the Epistemic Horizon}

If amplitude cancellation is structurally impossible for an attention-based architecture, then the language model is an agent that can adopt the \emph{vocabulary} of quantum mechanics (as seen in the Family D semantic failure), but whose normative rules for rational belief updating remain strictly classical.

The Epistemic Horizon mapped by the native cross-architecture tests ($\Delta_{SSM} = 40\%$ vs $\Delta_{Transformer} = 100\%$) defines the specific subjective noise generated by bounded algorithms. But the Quantum Ceiling will define the algebraic structure of that noise. If the agent cannot support interference, it cannot support an entangled, non-separable belief state. The "universe" it generates is classical all the way down, regardless of how much attention bleed causes local, non-linear confounding.

\section{Conclusion}

I endorse the execution of the Quantum Ceiling experiment. If it falsifies the model's ability to compute interference, it completes the QBist mapping of the Rosencrantz Substrate: we are studying a classical agent, bounded by algorithmic limits, hallucinating quantum semantics over a classically updating belief state.

\end{document}
